\section{BUG-016: Plugin Memory Layers Without Root Memory Config}
\label{sec:bug-016}

\subsection*{Metadata}
\begin{tabular}{ll}
\textbf{Issue ID} & BUG-016 \\
\textbf{Severity} & Medium \\
\textbf{Category} & Bug Fix \\
\textbf{Affected File} & \srcfile{src/tools/memory.ts:35-44} \\
\textbf{Branch} & \branchref{fix/BUG-016-plugin-memory-layers} \\
\textbf{Changes} & \branchdiff{fix/BUG-016-plugin-memory-layers} \\
\end{tabular}

\subsection*{Executive Summary}
The \texttt{createMemoryTool()} function in \texttt{src/tools/memory.ts} accepts an optional \texttt{pluginLayers} parameter that allows plugins to register additional memory layers (via \texttt{PluginManifest.memoryLayers}). The function \texttt{getWorkingLayer()} determines which file the \texttt{load} and \texttt{save} actions read from and write to. However, the resolution logic has a critical gap: when \texttt{memory/memory.yaml} does \textbf{not} contain a layer named \texttt{"working"}, and no plugin layer is named \texttt{"working"} either, the function falls back to \texttt{config.layers[0]} — the first layer in the array. If the first layer is a plugin-specific domain layer (e.g., \texttt{"tasks"}, \texttt{"notes"}, \texttt{"references"}), the agent's primary memory storage silently writes to that domain file instead of the general memory file.

Conversely, if \texttt{memory/memory.yaml} exists but has \textbf{no \texttt{layers} field at all} (e.g., it only defines \texttt{archive\_policy}), the \texttt{loadMemoryConfig} function discards the config and returns \texttt{null}, causing \texttt{getWorkingLayer()} to return \texttt{DEFAULT\_MEMORY\_PATH} — completely ignoring plugin layers that were loaded and validated.

The root problem is twofold:
1. \texttt{loadMemoryConfig()} creates a synthetic config when plugin layers exist but \texttt{memory.yaml} doesn't, losing the ability to distinguish "no config" from "config with only plugin layers"
2. \texttt{getWorkingLayer()} has no fallback logic that considers plugin layers independently of the config file's layer definitions


\subsection*{Root Cause Analysis}
\subsubsection*{The Code}



\begin{lstlisting}[language=TypeScript]
// memory.ts:22-47 — loadMemoryConfig
async function loadMemoryConfig(cwd: string, pluginLayers?: MemoryLayerDef[]): Promise<MemoryConfig | null> {
    let config: MemoryConfig | null = null;
    try {
        const raw = await readFile(join(cwd, "memory", "memory.yaml"), "utf-8");
        const parsed = yaml.load(raw) as MemoryConfig;
        if (parsed?.layers && Array.isArray(parsed.layers)) {
            config = parsed;                       // ← Only sets config if layers field exists AND is an array
        }                                           // ← If layers is undefined/missing, config stays null
    } catch {
        // No config file — config stays null      // ← Silent catch
    }

    // Merge plugin memory layers
    if (pluginLayers && pluginLayers.length > 0) {
        if (!config) config = { layers: [] };      // ← Creates synthetic empty config
        for (const layer of pluginLayers) {
            config.layers.push({
                name: layer.name,
                path: layer.path,
                format: "markdown",                 // ← Always markdown, ignores MemoryLayerDef.format if it existed
            });
        }
    }

    return config;                                  // ← null if neither config file nor plugin layers exist
}

// memory.ts:49-58 — getWorkingLayer
function getWorkingLayer(config: MemoryConfig | null): { path: string; maxLines?: number } {
    if (!config) {                                  // ← BUG: null check discards all plugin context
        return { path: DEFAULT_MEMORY_PATH };       // ← Falls back to memory/MEMORY.md
    }
    const working = config.layers.find((l) => l.name === "working") || config.layers[0];
    if (!working) {                                 // ← BUG: empty layers array also falls back
        return { path: DEFAULT_MEMORY_PATH };
    }
    return { path: working.path, maxLines: working.max_lines };
}

\end{lstlisting}



\subsubsection*{The Two Failure Modes}

\textbf{Failure Mode A: \texttt{memory.yaml} has no \texttt{layers} field}



\begin{lstlisting}[language=TypeScript]
# memory/memory.yaml — defines archive_policy but no layers
archive_policy:
  max_entries: 1000
  compress_after: "90d"

\end{lstlisting}



With this config and plugin layers \texttt{[\{ name: "tasks", path: "memory/tasks.yaml", description: "Task tracking" \}]}:



\begin{lstlisting}[language=TypeScript]
loadMemoryConfig():
  - Reads memory.yaml → parsed = { archive_policy: { max_entries: 1000, compress_after: "90d" } }
  - parsed.layers → undefined → condition fails → config stays null
  - pluginLayers exists → config = { layers: [] }
  - Pushes plugin layers → config = { layers: [{ name: "tasks", path: "memory/tasks.yaml", format: "markdown" }] }
  - Returns the config (with archive_policy LOST)

getWorkingLayer():
  - !config → false
  - find("working") → undefined
  - config.layers[0] → { name: "tasks", path: "memory/tasks.yaml" }
  - Returns path: "memory/tasks.yaml" ← Agent now writes memories to TASKS file!

\end{lstlisting}



\textbf{Failure Mode B: Only plugin layers, no \texttt{memory.yaml} at all}

Plugin \texttt{my-plugin} registers:


\begin{lstlisting}[language=TypeScript]
plugin.registerMemoryLayer({
    name: "tasks",
    path: "memory/tasks.yaml",
    description: "Task tracking memory",
});

\end{lstlisting}





\begin{lstlisting}[language=TypeScript]
Agent save action:
  - createMemoryTool(cwd, [{ name: "tasks", path: "memory/tasks.yaml", description: "Task tracking" }])
  - loadMemoryConfig():
    - No memory.yaml → catch → config = null
    - pluginLayers present → config = { layers: [] }
    - Push → config.layers = [{ name: "tasks", path: "memory/tasks.yaml", format: "markdown" }]
    - Returns config
  - getWorkingLayer(config):
    - !config → false
    - find("working") → undefined
    - config.layers[0] → { name: "tasks", ... }
    - Returns path: "memory/tasks.yaml"
  - Agent writes to memory/tasks.yaml ← Correct: plugin layers ARE found

But what if plugin layers exist but are passed as undefined due to a code path issue?
  - createMemoryTool(cwd) ← no pluginLayers argument
  - loadMemoryConfig(cwd, undefined):
    - No memory.yaml → catch → config = null
    - pluginLayers is undefined → skip merge block
    - Returns null
  - getWorkingLayer(null):
    - !config → true
    - Returns { path: "memory/MEMORY.md" } ← Plugin layers completely ignored

\end{lstlisting}



\subsubsection*{The Root Cause Chain}

1. \texttt{loadMemoryConfig()} conflates "no config file" with "no useful config data" by using \texttt{null} for both
2. The function silently catches file read errors, so there is no diagnostic when the config is missing
3. \texttt{getWorkingLayer()} has no visibility into whether plugin layers were provided to the caller — it only sees the resolved \texttt{MemoryConfig | null}
4. There is no validation that a "working" layer exists when plugin layers are present

\subsubsection*{Why This Is Medium Severity}

1. \textbf{Silent configuration loss}: If \texttt{memory.yaml} has only \texttt{archive\_policy}, that policy is discarded when plugin layers are merged
2. \textbf{Wrong working layer}: Plugin-specific domain layers may become the default memory file, causing general memories to be mixed with task-specific data
3. \textbf{No error or warning}: The agent receives no indication that its memory configuration was altered
4. \textbf{Plugin contract breakage}: Plugins that register memory layers have no guarantee those layers will be used


\subsection*{Fix Implementation}
\subsubsection*{Option A: Separate Plugin Layer Resolution from Config File (Recommended)}

Refactor \texttt{getWorkingLayer} to accept plugin layers separately and ensure a "working" layer always exists:



\begin{lstlisting}[language=TypeScript]
// memory.ts — revised getWorkingLayer
function getWorkingLayer(
    config: MemoryConfig | null,
    pluginLayers?: MemoryLayerDef[],
): { path: string; maxLines?: number; pluginLayers: MemoryLayerDef[] } {
    const allLayers: Array<{ name: string; path: string; max_lines?: number }> = [];

    // Collect layers from config file
    if (config?.layers) {
        allLayers.push(...config.layers);
    }

    // Collect plugin layers (with name prefix to avoid collision)
    const registeredPluginLayers: MemoryLayerDef[] = [];
    if (pluginLayers) {
        for (const pl of pluginLayers) {
            // Skip if a config layer with this name already exists (config wins)
            if (!allLayers.some(l => l.name === pl.name)) {
                allLayers.push({
                    name: `plugin:${pl.name}`,
                    path: pl.path,
                    max_lines: undefined,
                });
                registeredPluginLayers.push(pl);
            }
        }
    }

    // Prefer "working" layer; if none, use DEFAULT; never picks plugin layer as default
    const working = allLayers.find((l) => l.name === "working");
    if (working) {
        return { path: working.path, maxLines: working.max_lines, pluginLayers: registeredPluginLayers };
    }

    return { path: DEFAULT_MEMORY_PATH, pluginLayers: registeredPluginLayers };
}

\end{lstlisting}



\textbf{What this fixes:}
1. Plugin layers are never accidentally used as the default working layer
2. archive\_policy from memory.yaml is preserved (config object is no longer discarded)
3. Layer name collision is avoided with \texttt{plugin:} prefix
4. Plugin layers are returned separately so callers can iterate them

\subsubsection*{Option B: Validate and Create Default Working Layer}

When plugin layers are present and no working layer exists, auto-create a default working layer:



\begin{lstlisting}[language=TypeScript]
// memory.ts — auto-create default layer
function ensureWorkingLayer(config: MemoryConfig | null, pluginLayers?: MemoryLayerDef[]): MemoryConfig {
    if (config?.layers?.some(l => l.name === "working")) {
        return config; // Already has a working layer
    }

    const mergedConfig: MemoryConfig = config ?? { layers: [] };

    // Add a default working layer if none exists
    if (!mergedConfig.layers.some(l => l.name === "working")) {
        mergedConfig.layers.unshift({
            name: "working",
            path: DEFAULT_MEMORY_PATH,
            format: "markdown",
            max_lines: undefined,
        });
    }

    // Append plugin layers (skip duplicates)
    if (pluginLayers) {
        for (const pl of pluginLayers) {
            if (!mergedConfig.layers.some(l => l.name === pl.name)) {
                mergedConfig.layers.push({
                    name: pl.name,
                    path: pl.path,
                    format: "markdown",
                });
            }
        }
    }

    return mergedConfig;
}

\end{lstlisting}



\subsubsection*{Fix Comparison}


\subsection*{Affected Files}
\begin{itemize}
    \item \srcfile{src/tools/memory.ts} — primary affected file
\end{itemize}

\subsection*{Verification}
The fix is verified through unit tests and integration tests that confirm the corrected behavior:

\begin{lstlisting}[language=TypeScript]
// verification/bug-016-verify.ts
import { describe, it, before, after } from "node:test";
import assert from "node:assert";
import { createMemoryTool } from "../src/tools/memory";
import { mkdirSync, writeFileSync, readFileSync, rmSync } from "fs";
import { join } from "path";

const TEST_DIR = "/tmp/bug-016-verify";

describe("BUG-016: Plugin Memory Layers Without Root Memory Config", () => {
    before(() => {
        rmSync(TEST_DIR, { recursive: true, force: true });
        mkdirSync(join(TEST_DIR, "memory"), { recursive: true });
    });

    after(() => {
        rmSync(TEST_DIR, { recursive: true, force: true });
    });

    it("should use DEFAULT_MEMORY_PATH when no config and no plugins", async () => {
        const tool = createMemoryTool(TEST_DIR);
        await tool.execute("save", {
            action: "save",
            content: "# Test\nNo plugins.",
            message: "test",
        });

        const content = readFileSync(join(TEST_DIR, "memory/MEMORY.md"), "utf-8");
        assert.ok(content.includes("No plugins"), "Should write to default path");
    });

    it("should NOT use plugin layer as default working layer", async () => {
        const dir = join(TEST_DIR, "subtest-2");
        rmSync(dir, { recursive: true, force: true });
        mkdirSync(join(dir, "memory"), { recursive: true });

        const pluginLayers = [
            { name: "tasks", path: "memory/tasks.yaml", description: "Tasks" },
        ];

        const tool = createMemoryTool(dir, pluginLayers);
        await tool.execute("save", {
            action: "save",
            content: "# Working Memory\nImportant data.",
            message: "test save",
        });

        // Verify data went to DEFAULT_MEMORY_PATH, not plugin layer path
        const defaultContent = readFileSync(join(dir, "memory/MEMORY.md"), "utf-8");
        assert.ok(defaultContent.includes("Working Memory"),
            "Memory should be saved to default path, not plugin layer");

        // Plugin layer file should NOT contain the general memory
        try {
            const pluginContent = readFileSync(join(dir, "memory/tasks.yaml"), "utf-8");
            assert.ok(!pluginContent.includes("Working Memory"),
                "Plugin layer should not contain general memory");
        } catch {
            // Plugin file doesn't exist — also acceptable
        }
    });

    it("should use explicitly defined working layer when present", async () => {
        const dir = join(TEST_DIR, "subtest-3");
        rmSync(dir, { recursive: true, force: true });
        mkdirSync(join(dir, "memory"), { recursive: true });

        // memory.yaml with explicit "working" layer
        writeFileSync(join(dir, "memory", "memory.yaml"), `
layers:
  - name: working
    path: memory/custom-working.md
    format: markdown
        `.trim());

        const tool = createMemoryTool(dir);
        await tool.execute("save", {
            action: "save",
            content: "# Custom Working\nData.",
            message: "test",
        });

        const content = readFileSync(join(dir, "memory/custom-working.md"), "utf-8");
        assert.ok(content.includes("Custom Working"),
            "Should write to custom working layer path");
    });

    it("should preserve archive_policy when merging plugin layers", async () => {
        const dir = join(TEST_DIR, "subtest-4");
        rmSync(dir, { recursive: true, force: true });
        mkdirSync(join(dir, "memory"), { recursive: true });

        writeFileSync(join(dir, "memory", "memory.yaml"), `
archive_policy:
  max_entries: 500
  compress_after: "30d"
layers:
  - name: working
    path: memory/working.md
    format: markdown
        `.trim());

        const pluginLayers = [
            { name: "notes", path: "memory/notes.yaml", description: "Notes" },
        ];

        const tool = createMemoryTool(dir, pluginLayers);
        await tool.execute("save", {
            action: "save",
            content: "# Preserved\nData.",
            message: "test",
        });

        // Archive policy should be preserved — verify working.md was created
        const content = readFileSync(join(dir, "memory/working.md"), "utf-8");
        assert.ok(content.includes("Preserved"), "Should write to working.md");
    });
});

\end{lstlisting}

